% =====================================================================
%  Rebuttal to Reviewer Comments
%  Energy-Efficient and Carbon-Aware Scheduling in Cloud Computing
%  using an Enhanced Whale Optimization Algorithm (CA-WOA)
%  Navaneetha Thalakokkula (23213914) - MSc Cloud Computing, NCI
%
%  Items marked [TO BE COMPLETED BY SUPERVISOR] are left blank
%  deliberately: they are decisions or judgements that are not mine to
%  make. Every other response points to a specific section, table or
%  figure in the revised manuscript.
% =====================================================================
\documentclass[11pt,a4paper]{article}
\usepackage[utf8]{inputenc}
\usepackage[T1]{fontenc}
\usepackage{mathptmx}
\usepackage[margin=1in]{geometry}
\usepackage{amsmath,amssymb}
\usepackage{xcolor}
\usepackage{enumitem}
\usepackage{booktabs}
\usepackage[hidelinks]{hyperref}
\usepackage{tcolorbox}

\definecolor{cmt}{rgb}{0.30,0.30,0.30}
\newcommand{\rsp}[1]{\vspace{2pt}\noindent\textbf{Response.} #1\par\vspace{6pt}}
\newcommand{\rwhere}[1]{\noindent\textit{Where addressed:} #1\par\vspace{10pt}}
\newcommand{\rpending}{\vspace{2pt}\noindent\textbf{Response.}
  \textcolor{red}{\textbf{[TO BE COMPLETED BY SUPERVISOR]}}\par\vspace{10pt}}
\newenvironment{revcomment}[1]
  {\begin{tcolorbox}[colback=gray!6,colframe=gray!45,boxrule=0.4pt,
     left=6pt,right=6pt,top=4pt,bottom=4pt]\textbf{#1}\\[2pt]\itshape\color{cmt}}
  {\end{tcolorbox}\vspace{4pt}}

\title{\textbf{Rebuttal to Reviewer Comments}\\[4pt]
\large Energy-Efficient and Carbon-Aware Scheduling in Cloud Computing\\
using an Enhanced Whale Optimization Algorithm (CA-WOA)}
\author{Navaneetha Thalakokkula\\Student ID: 23213914\\
MSc Cloud Computing, National College of Ireland\\Supervisor: Punit Gupta}
\date{\today}

\begin{document}
\maketitle

\section*{Summary of the revision}
We thank the reviewers for a detailed and constructive set of comments. All 43
points have been addressed. The revision is substantial: the experimental scale
increased from 60 tasks and 3--5 seeds to 500--3000 tasks on 10 and 20 hosts with
30 independent seeds, and from one carbon-intensity trace to five real grid
regions. Twenty-nine experiments are reported, releasing one row per
configuration, method and seed.

Four changes deserve to be stated at the outset, because they correct or withdraw
claims made in the submitted version rather than merely extending it.

\begin{enumerate}[leftmargin=*]
\item \textbf{The novelty claim has been withdrawn.} The submitted abstract stated
that WOA ``has not previously been used for carbon''. We no longer make any
priority claim. The contribution is restated as carbon-aware \emph{seeding}, an
initialisation strategy demonstrated to transfer across six optimisers.

\item \textbf{The equal-budget claim was false and is corrected.} Section~6.1 of
the submitted version asserted identical evaluation budgets. Measured over all
runs, HHO consumes 8{,}908 objective evaluations against 4{,}840 for every other
method---roughly $1.84\times$. This is now stated explicitly and the evaluation
count is recorded for every run.

\item \textbf{The GA baseline performed no search.} \texttt{GA.OriginalGA} in
Mealpy 3.0.3 evaluates only its initial population: 40 objective evaluations
rather than 4{,}840. The submitted GA results were therefore not a genetic
algorithm. \texttt{GA.BaseGA} is used throughout the revision and the defect is
quantified rather than silently fixed.

\item \textbf{The forecasting headline did not reproduce.} The reported 12\,\%
improvement over a baseline LSTM becomes approximately 7\,\% once TensorFlow
operation-level determinism is enabled and results are averaged over five
repeats, and the CNN-LSTM component is in fact the \emph{weakest} learned model at
every horizon tested. The claim is withdrawn and replaced with the 12-hour
horizon result, which is defensible.
\end{enumerate}

\noindent Two findings that \emph{limit} the contribution are also reported rather
than omitted: under an enforced minimum active-host constraint the advantage of
seeding over standard WOA falls to $+0.04$ percentage points, and on a grid whose
carbon intensity varies only $1.78\times$ across the day the proposed method loses
to the best competing method. Both are in Section~6.17 (Limitations).

All artefacts are archived at \url{https://doi.org/10.5281/zenodo.22125797}.

\bigskip
\hrule
\bigskip

\section{Reviewer 2}

\subsection*{Host and power model}
\begin{revcomment}{Comment R2.1}
Justify $P_{\text{idle}} = 100$\,W, $P_{\text{max}} = 250$\,W and normalized
capacity $C = 1$. Consider heterogeneous hosts, non-linear power models,
startup/shutdown overheads and minimum active-host constraints.
\end{revcomment}
\rsp{The endpoints are now justified in the text as a mid-range two-socket volume
server: idle draw of roughly 40\,\% of peak is typical of published SPECpower
results, and the 150\,W dynamic range spans the CPU-dominated portion of the
envelope. Because a linear model flatters consolidation, two further models with
identical endpoints were added: a convex \emph{cubic} model and a concave
\emph{piecewise} model interpolating a measured SPECpower-style curve. CA-WOA
ranks first in all six configurations under each of the three models---18 of 18---so
the conclusions are not an artefact of the linear assumption.

Heterogeneous fleets were added: three fleets (uniform; \emph{mixed}, where newer
hosts are larger and more efficient; and \emph{skewed}, deliberately adversarial,
where the large hosts are the inefficient ones), with hosts filled in order of
energy per unit of work delivered. Startup energy is charged per off-to-on
transition; the greedy schedule at 1500 tasks triggers 58 host boots, so this is a
material cost rather than a rounding term. Seeding helps in 52 of 54 cells, and its
advantage is \emph{larger} on heterogeneous fleets ($+2.96$ against $+1.55$
percentage points), since packing onto unequal hosts is a harder problem.

A workload-determined minimum active-host constraint was added, modelling a warm
pool that cannot scale to zero. This reverses the comparison with the greedy
heuristic in our favour---12 of 12 configurations rather than 3 of 12---because
deferral now costs idle energy. It also removes most of the benefit of seeding,
which we report as a limitation.}
\rwhere{Section~3.2 (energy and carbon model); Section~6.9 (power models); Section~6.10
(minimum active-host constraint); Section~6.11 (heterogeneous fleets and startup
overheads).}

\subsection*{Novelty claim}
\begin{revcomment}{Comment R2.2}
Strengthen the literature review and moderate unsupported claims such as ``first
application.'' Clearly identify the contribution of carbon-aware seeding,
SLA-aware optimization and predictive forecasting.
\end{revcomment}
\rsp{The priority claim has been removed from the abstract, the introduction and
the research-gap subsection. We now state plainly that WOA has been applied to
cloud task scheduling in enhanced form, and that carbon-aware metaheuristic
scheduling is established, citing work from 2022--2026. Twelve references have been
added, every DOI resolved against Crossref before use, and Table~1 now spans
2022--2026 rather than stopping at 2023.

The contribution is restated as five specific items, of which the central one is
carbon-aware seeding: initialising part of a population-based optimiser with a
carbon-aware greedy schedule. It improves final solution quality in 59 of 60
(optimiser, configuration) cells across WOA, GWO, PSO, DE, HHO and GA, and
Corollary~2 proves it can never be worse than the schedule it seeds. Because it
transfers across optimisers, it is a property of the initialisation rather than of
WOA.

The headline figure is also decomposed, since most of it is not attributable to
the proposed method: of the 83.25\,\% reduction, 80.50 percentage points come from
consolidation alone, temporal shifting adds 2.06, and CA-WOA adds 0.69 over the
carbon-aware greedy heuristic.}
\rwhere{Abstract; Section~1 (contributions); Section~2.3; Section~2.5;
Table~1; Section~6.3 (host-capacity model comparison).}

\subsection*{Ablation study}
\begin{revcomment}{Comment R2.3}
Separate the contribution of standard WOA, improved/random initialization,
carbon-aware initialization, and the forecasting component.
\end{revcomment}
\rsp{A five-arm ablation is reported under both energy models, with 30 seeds. Two
results carry the argument. First, carbon-aware seeding improves quality in 59 of
60 cells (mean $+2.51$ points under the published model, significant in 36 of 36;
$+0.91$ under the capacity constraint, significant in 21 of 24). Second---and this
is the result that isolates \emph{what} the seeding contributes---an ``improved''
initialisation using Latin-hypercube plus opposition-based sampling but \emph{no
carbon data} is statistically indistinguishable from uniform random ($+0.003$ to
$+0.027$ points, $p$ = 0.06 to 0.94). The benefit therefore comes from the domain
information, not from a more sophisticated sampling scheme.

The forecasting component is ablated separately as predicted versus reactive
versus perfect-foresight scheduling at identical optimisation budgets.}
\rwhere{Section~6.7 (ablation) and Table~7; Section~6.15 (forecasting).}

\subsection*{Fitness-weight sensitivity}
\begin{revcomment}{Comment R2.4}
Evaluate different $\alpha$, $\beta$ and $\gamma$ values and justify the selected
weights $\alpha = 0.4$, $\beta = 0.3$ and $\gamma = 0.3$.
\end{revcomment}
\rsp{Eleven weight combinations were evaluated over 30 seeds. Read on carbon alone
the result is misleading: $\alpha = 1.0$ attains the highest carbon reduction in
the study (89.92\,\%) but violates 71.1\,\% of deadlines and overloads hosts by
33.5\,\%. Removing $\beta$ or $\gamma$ likewise buys carbon only by breaking the
constraint that term exists to protect. Among settings that hold both SLA
violations and overload near zero, the chosen $(0.4, 0.3, 0.3)$ sits on the knee of
the trade-off at 87.58\,\% carbon with 0.13\,\% SLA violations. The weights are
therefore justified by the constraint surface rather than by tuning for the
headline number.

The seeded fraction was also swept: carbon varies by less than 0.3 points across
fractions from 1/6 to 1, but variance falls sharply once seeding is enabled
(standard deviation 0.381 to 0.111), so 1/3 is a reasonable rather than a uniquely
optimal choice.}
\rwhere{Section~6.8 (fitness-weight sensitivity); Table~7 (rows with $\beta=0$ and $\gamma=0$).}

\subsection*{SLA validation}
\begin{revcomment}{Comment R2.5}
Provide absolute violation counts, per-run results and more demanding workloads.
Consider hard constraints or a repair mechanism if zero SLA violation is a central
claim.
\end{revcomment}
\rsp{Absolute violation counts are now reported alongside percentages, and raw
results are released with one row per configuration, method and seed, so per-run
values are directly inspectable. Workloads increased from 60 tasks to 500--3000.

A repair decoding was implemented, which clamps every start time to the task's
deadline-feasible window, making deadlines hard by construction. Under repair
decoding the maximum SLA violation across the entire experiment is exactly
0.0000\,\%. This replaces the earlier reliance on a penalty term. For context, the
aggressive optimisers under soft constraints violate a mean of 58.0\,\% (PSO),
50.7\,\% (DE) and 35.4\,\% (GA) of deadlines at scale, so they are worse on both
axes simultaneously.}
\rwhere{Section~3.3 (metrics); Section~6.5 (capacity as a feasibility constraint); Table~6.}

\subsection*{Statistical validation}
\begin{revcomment}{Comment R2.6}
Use sufficient independent runs and report confidence intervals, significance
tests, effect sizes and preferably per-seed results.
\end{revcomment}
\rsp{Every algorithm is now run over 30 independent seeds. Each comparison reports
mean, standard deviation and a 95\,\% confidence interval. Significance is assessed
with \emph{both} Welch's $t$-test and the Mann--Whitney $U$ test---several methods
have near-zero variance, where a $t$-test is ill-conditioned---and a difference is
reported as significant only when both agree at $\alpha = 0.05$. Effect sizes are
given as Cohen's $d$ and Cliff's $\delta$. Per-seed results are released in full.}
\rwhere{Section~3.3; all result tables; archived per-seed CSV files.}

\subsection*{Forecasting validation}
\begin{revcomment}{Comment R2.7}
Report forecast horizon, prediction procedure, additional metrics such as RMSE and
MAPE/MASE where appropriate, and suitable naive/seasonal baselines.
\end{revcomment}
\rsp{The forecaster is now evaluated with direct multi-step prediction at horizons of
1, 6, 12, 24 and 48 slots, using MAE, RMSE, MAPE and MASE, against three naive
baselines: persistence, seasonal-naive at 24 hours, and seasonal-naive at one week.
Each model is trained five times and results are reported as mean $\pm$ standard
deviation, because the single-run figures in the submitted version did not
reproduce.

At the 12-hour horizon the ensemble achieves MAE 26.09 $\pm$ 0.48\,gCO$_2$/kWh
against 37.03 for the strongest naive baseline, a MASE of 0.59. Full
hyperparameters are now tabulated.}
\rwhere{Section~3.3; Section~6.15 (forecasting); Table~9; Figure~8.}

\subsection*{Forecast--scheduling coupling}
\begin{revcomment}{Comment R2.8}
Clearly describe how predictions are supplied to CA-WOA, including forecast
horizon, update frequency, online operation and the effect of forecast errors.
\end{revcomment}
\rsp{The coupling is now described explicitly and measured. A 12-hour-ahead
ensemble forecast replaces the true carbon signal in the objective, and the
resulting schedule is scored against the \emph{actual} signal. Predicted,
reactive and perfect-foresight scheduling are compared at identical optimisation
budgets. The forecast captures 73--89\,\% of the advantage that perfect foresight
holds over purely reactive scheduling.

We note that an earlier version of our own analysis reported this as 99--100\,\%.
That figure divided the two total carbon reductions, both of which are dominated
by the consolidation baseline, and would therefore have flattered any forecast
regardless of accuracy. The corrected quantity is the fraction of the
\emph{advantage over reactive scheduling} that the forecast recovers.

Forecast-error degradation is measured separately by injecting controlled error and
observing the resulting scheduling penalty.}
\rwhere{Section~6.15 (forecasting); Section~3.3.}

\subsection*{Computational scalability}
\begin{revcomment}{Comment R2.9}
Report runtime, memory usage, total objective-function evaluations and convergence
behaviour as task count increases.
\end{revcomment}
\rsp{Runtime, peak memory and the number of function evaluations are recorded for
every run, and convergence curves are retained per epoch. Runtime scales linearly
in task count: WOA is the cheapest search at 1.54\,s and 0.99\,MB at 500 tasks,
rising to 6.20\,s and 4.87\,MB at 3000; GA is the most expensive at 2.61 to
10.57\,s. Peak memory never exceeded 9.6\,MB in any run.

Convergence behaviour shows the seeding acting as designed: CA-WOA begins at
roughly half the fitness of standard WOA and stabilises within 1\,\% of its final
value by epoch 8 against 16, reaching a better final value. We also report a
caveat that favours the competitors: GA and GWO had \emph{not} converged at the
120-epoch budget, requiring 102 and 112 epochs respectively, so their results are
lower bounds and this comparison understates them.}
\rwhere{Section~6.14 (computational cost and convergence); Figure~7.}

\subsection*{Reproducibility}
\begin{revcomment}{Comment R2.10}
Provide precise dataset references, preprocessing details, configuration files,
seeds, code repository/persistent link or DOI, and sufficient artefacts to
reproduce the experiments.
\end{revcomment}
\rsp{All code, data, per-seed raw results and the analysis scripts are archived at
\url{https://doi.org/10.5281/zenodo.22125797} and mirrored on GitHub. Dataset
references and preprocessing are specified in Section~3.1, including the sampling
correction described under Comment~R2.11 below. Seeds are fixed and enumerated
(1--30); dependency versions are recorded.

The artefact includes a 27-check validation gate that must pass before any new
result is trusted. It confirms that the fast vectorised evaluator is numerically
identical to the original implementation (worst relative deviation
$9.3 \times 10^{-16}$, machine precision) and that the setup reproduces the
originally published Table~4 exactly. Two experiments are deliberately retained
and labelled \texttt{INVALID}, documenting modelling errors that we detected and
corrected during the revision.}
\rwhere{Section~3.3; Availability of Data and Materials; archived artefact.}

\bigskip
\hrule
\bigskip

\section{Reviewer 2: Essential additional experiments}

\begin{revcomment}{Comment R2.11}
Ablation: WOA; WOA + carbon-aware seeding; CA-WOA without SLA term; CA-WOA without
overload/cost term; full CA-WOA.
\end{revcomment}
\rsp{All five arms are reported, plus the greedy-only baseline and the carbon-blind
``improved'' initialisation, over 30 seeds under both energy models. See
Comment~R2.3.

We also report a data-preparation defect found during the revision, because it
affects the interpretation of the submitted results. The submitted workload was
drawn with \texttt{head(60)}, which selected tasks sharing a single submit
timestamp: every arrival collapsed to slot 0 and the workload occupied 4 of the 144
available slots, leaving a temporal-shifting scheduler almost nothing to shift.
The same concentration inflated the derived host count from 1 to 5. Tasks are now
sampled at an even stride across all 52{,}057 usable records, preserving the real
submission ordering.}
\rwhere{Section~6.7 (ablation); Section~3.1 (data sources).}

\begin{revcomment}{Comment R2.12}
Compare against an earliest-deadline-first + lowest-carbon-slot heuristic and a
constrained-scheduling optimization method.
\end{revcomment}
\rsp{The EDF + lowest-carbon-slot heuristic is implemented and appears in every
results table. It is a strong baseline in the unconstrained regime---so strong that
we prove it is \emph{optimal} there (Lemma~1): in the temporal model the objective
is separable across tasks, so choosing each task's lowest-carbon feasible slot
independently attains the global optimum and no search can improve on it. This
explains the otherwise puzzling exact tie between CA-WOA and the greedy baseline
in the submitted results.

Under host-capacity constraints, however, the same heuristic becomes infeasible: it
piles work into the greenest slots and overloads hosts by up to 45.23\,\%. Across
five real grid regions it violates capacity in 58--75\,\% of three-day windows.
Its apparently superior carbon figures are therefore an infeasible upper bound
rather than a competing result.

Following our supervisor's direction, the constrained-scheduling optimisation
method is addressed \emph{theoretically} rather than implemented: the problem is
formulated exactly as a mixed-integer linear program, with the complexity argument
and the role of an exact solve as an optimality-gap reference given in
Section~4.4. We identify a full MILP solve as future work.}
\rwhere{Section~4.4 (exact formulation, Lemma~1); Section~6.5; all results tables.}

\begin{revcomment}{Comment R2.13}
Sensitivity analysis for fitness weights and seeded-population percentage.
\end{revcomment}
\rsp{Both are reported: eleven weight combinations and six seeded fractions from 0 to
1, each over 30 seeds. See Comment~R2.4.}
\rwhere{Section~6.8 (fitness-weight sensitivity).}

\begin{revcomment}{Comment R2.14}
Use at least 20--30 independent random seeds, or provide a statistically justified
alternative with confidence intervals.
\end{revcomment}
\rsp{Thirty independent seeds are used throughout, with 95\,\% confidence intervals,
two significance tests and two effect-size measures. See Comment~R2.6.}
\rwhere{Section~3.3; all results tables.}

\begin{revcomment}{Comment R2.15}
Report convergence curves and total objective-function evaluations.
\end{revcomment}
\rsp{Convergence curves are reported per epoch with dispersion bands, and the
evaluation count is recorded for every run---which is how the unequal-budget
problem was detected. See Comment~R2.9.}
\rwhere{Figure~7; Section~6.14 (computational cost and convergence).}

\begin{revcomment}{Comment R2.16}
Report runtime and memory consumption for all algorithms at 50, 100, 200 and 300
tasks.
\end{revcomment}
\rsp{Runtime and memory are reported at exactly those four task counts, in addition
to the 500--3000 range used for the main study. WOA runs in 0.89\,s at 50 tasks and
1.24\,s at 300.}
\rwhere{Section~6.14 (computational cost and convergence).}

\begin{revcomment}{Comment R2.17}
Evaluate multiple workload subsets rather than relying on one fixed 60-task sample.
\end{revcomment}
\rsp{Four disjoint subsets of the Google trace are evaluated, and the NASA-iPSC log
from the Parallel Workloads Archive is added as an independent workload with a
genuine 92-day arrival process. CA-WOA is stable across the four subsets---standard
deviation 0.13 to 0.59 percentage points depending on configuration---and wins all six
NASA configurations.

Absolute carbon reductions on NASA are markedly lower, 61.6 to 67.6\,\% against 84 to
89\,\% on the Google trace, because a real 92-day arrival process offers far less
opportunity to shift work into a low-carbon slot. We report that as the more realistic
figure rather than the higher one. The SLA contrast, however, widens: averaged per
configuration, PSO misses 52.5 to 63.6\,\% of deadlines, DE 46.3 to 56.9\,\% and GA
21.2 to 49.9\,\%, while CA-WOA stays between 0.00 and 0.41\,\%.}
\rwhere{Section~6.13 (multiple workload subsets and a second trace); Table~8;
Section~3.1 (data sources).}

\begin{revcomment}{Comment R2.18}
Test heterogeneous host capacities and more realistic power-consumption models.
\end{revcomment}
\rsp{Both are now tested: three fleet compositions including an adversarial one, and
three power models. See Comment~R2.1.}
\rwhere{Sections~6.9 (power models) and 6.11 (heterogeneous fleets).}

\begin{revcomment}{Comment R2.19}
Evaluate multiple carbon-intensity traces/regions or justify the representativeness
of the UK signal.
\end{revcomment}
\rsp{Five real half-hourly regional traces are now used, spanning the same 111-day
period: the United Kingdom (National Grid ESO), the Republic of Ireland and
Northern Ireland (EirGrid Smart Grid Dashboard), and California and the
Mid-Atlantic (US EIA, with intensity derived from the hourly generation mix using
IPCC AR5 median lifecycle emission factors). They differ substantially, with mean
intensity from 111.9 to 333.5\,gCO$_2$/kWh and diurnal variability from
$12.30\times$ down to $1.78\times$.

Separately, the experiment was repeated across all 37 disjoint three-day windows in
the UK record, confirming that the single window used in the submitted version was
not unrepresentative.}
\rwhere{Section~3.1 and Table~2; Section~6.12 (multiple carbon regions); Figure~6.}

\begin{revcomment}{Comment R2.20}
Test forecast degradation and quantify scheduling performance under increasing
forecast error.
\end{revcomment}
\rsp{Controlled error is injected into the carbon signal at increasing magnitudes and
the resulting scheduling penalty is measured, so the sensitivity of carbon
reduction to forecast quality is quantified directly rather than inferred.}
\rwhere{Section~6.15 (forecasting).}

\begin{revcomment}{Comment R2.21}
Compare predicted-carbon, reactive and perfect-foresight scheduling using identical
optimization budgets.
\end{revcomment}
\rsp{All three are compared at an identical budget of 4{,}840 objective evaluations.
The forecast recovers 73--89\,\% of the advantage that perfect foresight holds over
reactive scheduling. See Comment~R2.8.}
\rwhere{Section~6.15 (forecasting).}

\bigskip
\hrule
\bigskip

\section{Reviewer 2: Presentation corrections}

\begin{revcomment}{Comment R2.22}
Correct grammatical and wording errors.
\end{revcomment}
\rsp{The manuscript has been revised throughout. Specific instances addressed
include vague quantifiers (``up to about 87\,\%'' replaced with the figure and its
configuration), unsupported editorialising about the proposed method, and undefined
qualitative claims such as ``scales safely'', which is now stated as the
measurable property it refers to.}
\rwhere{Throughout.}

\begin{revcomment}{Comment R2.23}
Use consistent terminology for carbon intensity, carbon-aware scheduling, carbon
reduction and emissions.
\end{revcomment}
\rsp{One term is now used per concept throughout: \emph{carbon intensity}
(gCO$_2$/kWh), \emph{carbon-aware scheduling}, \emph{carbon reduction} (\%) and
\emph{emissions} (kgCO$_2$).

One inconsistency deserves specific mention. The submitted version used ``host''
and ``VM'' interchangeably, including ``10 VMs and 20 VMs''. The simulation places
\emph{tasks} on \emph{hosts} of capacity $C$ and does not model virtual machines at
all. The term VM has been removed and a note added making the distinction
explicit.}
\rwhere{Abbreviations block; throughout.}

\begin{revcomment}{Comment R2.24}
Define all abbreviations at first use.
\end{revcomment}
\rsp{An abbreviations block has been added after the keywords, defining nineteen
terms. It also records that ``CI'' is reserved for carbon intensity, with
``95\,\% confidence interval'' always written out in full, since the submitted
version used CI for both.}
\rwhere{Abbreviations block, after the keywords.}

\begin{revcomment}{Comment R2.25}
Improve figure readability and include units, legends and error information.
\end{revcomment}
\rsp{All figures have been regenerated from the raw per-seed results at 300\,dpi.
Every axis label carries units; legends sit clear of the data; every mean carries a
95\,\% confidence interval or standard deviation; and the proposed method is
distinguished by marker and hatch as well as colour so the figures survive
greyscale printing. Two new figures have been added, showing convergence behaviour
and the relationship between seeding benefit and grid carbon variability.

Where a figure uses a truncated vertical axis, the caption states the actual spread
so the visual separation is not read as larger than it is.}
\rwhere{Figures 3--8 (all regenerated); Figures~7 and~8 are new.}

\begin{revcomment}{Comment R2.26}
Add standard deviations or confidence intervals to relevant tables.
\end{revcomment}
\rsp{All result tables now report dispersion over 30 seeds. See Comment~R2.6.}
\rwhere{All results tables.}

\begin{revcomment}{Comment R2.27}
Provide complete forecasting-model hyperparameters.
\end{revcomment}
\rsp{Full hyperparameters are tabulated: look-back window, forecast horizons,
convolution filters and kernel size, pooling, recurrent unit counts, gradient
boosting configuration, optimiser, loss, epochs, batch size, train/test split,
feature set, scaling procedure (fitted on the training split only) and the random
seed.}
\rwhere{Table~9; Section~6.15 (forecasting).}

\begin{revcomment}{Comment R2.28}
Correct the ``Availability of data and materials: NA'' statement because public
datasets are used.
\end{revcomment}
\rsp{Corrected. A full availability statement now names all three workload and five
carbon data sources, all of which are public, together with the archived artefact
DOI.}
\rwhere{Availability of Data and Materials.}

\begin{revcomment}{Comment R2.29}
Provide a persistent code repository link/DOI if permitted.
\end{revcomment}
\rsp{The artefact is archived with a persistent DOI:
\url{https://doi.org/10.5281/zenodo.22125797}. The concept DOI is cited so that it
continues to resolve if further versions are deposited.}
\rwhere{Availability of Data and Materials.}

\begin{revcomment}{Comment R2.30}
Clarify whether time-of-use tariff values are actually included in the reported
objective/results.
\end{revcomment}
\rsp{Clarified, and the honest answer is that the tariff is present but carries
almost no information. It enters the objective through the cost term, but because
every method schedules the same total workload and the tariff peak does not
coincide with the carbon peak, cost is effectively constant: identical to three
decimal places at 0.755\,GBP for every method in the temporal-model configuration,
and varying only between 0.796 and 0.935\,GBP in the capacity model, without
altering any ranking. Cost is therefore reported for completeness and explicitly
\emph{not} presented as a co-optimised objective. We also state that the UK tariff
is applied unchanged to the non-UK regions as a cost proxy, not as a claim about
local pricing.}
\rwhere{Section~3.2.}

\begin{revcomment}{Comment R2.31}
Use consistent units for energy, emissions, carbon intensity and monetary cost.
\end{revcomment}
\rsp{Units are now consistent throughout and stated in every table header: energy in
kWh, emissions in kgCO$_2$, carbon intensity in gCO$_2$/kWh, power in W, cost in
GBP, makespan in hours, runtime in seconds and memory in MB.}
\rwhere{Throughout; all table headers.}

\begin{revcomment}{Comment R2.32}
Check reference metadata, DOI information and formatting.
\end{revcomment}
\rsp{Every entry in the bibliography was resolved against Crossref before use. The
reference list has grown from 19 entries to 31, all cited in the text, all carrying
DOIs. Two metadata inconsistencies in the original list were corrected.

One entry was deliberately \emph{not} added: a preprint for which only an SSRN
record exists and the journal version is still in press. We prefer to omit it than
to cite a version of record that does not yet exist.}
\rwhere{Bibliography.}

\bigskip
\hrule
\bigskip

\section{Reviewer 3}

\begin{revcomment}{Comment R3.1}
Add evaluation on multiple carbon intensity datasets and newer traces.
\end{revcomment}
\rsp{Five real regional carbon-intensity datasets are now used, and the NASA-iPSC
workload trace has been added alongside the Google trace. See Comment~R2.19.}
\rwhere{Section~3.1, Table~2; Section~6.12 (multiple carbon regions).}

\begin{revcomment}{Comment R3.2}
Include stronger baselines (at least one carbon-aware heuristic, one modern
scheduler).
\end{revcomment}
\rsp{Two published policies were implemented in the same simulator, using no fitness
function and no metaheuristic machinery, so they are independent of the proposed
method: the virtual-capacity-curve load-shaping policy of Google's production
carbon-aware scheduler, and threshold-based temporal shifting from
\emph{Let's Wait Awhile}. The VCC throttle floor was swept over eight values so the
baseline is not weakened by an arbitrary parameter choice; it proved insensitive.

Only the virtual-capacity-curve policy and the population-based methods are
feasible in every configuration. Among these, CA-WOA wins all twelve
configurations, by 4.48 percentage points at 500 tasks narrowing to 0.35 at 3000.
The threshold policy shifts load toward greener slots but has no capacity
awareness and remains infeasible, which shows that carbon awareness alone is
insufficient.}
\rwhere{Section~6.6 (comparison with published schedulers); Table~6.}

\begin{revcomment}{Comment R3.3}
Provide statistical significance tests for all comparisons.
\end{revcomment}
\rsp{Both Welch's $t$-test and the Mann--Whitney $U$ test are applied to every
comparison, with significance claimed only when both agree at $\alpha = 0.05$,
alongside Cohen's $d$ and Cliff's $\delta$. See Comment~R2.6.}
\rwhere{Section~3.3; all results tables.}

\begin{revcomment}{Comment R3.4}
Clarify the forecast horizon and include multi-step forecasting evaluation.
\end{revcomment}
\rsp{Direct multi-step prediction is now evaluated at five horizons from 30 minutes
to 24 hours, with four error metrics and three naive baselines. See
Comment~R2.7.}
\rwhere{Section~6.15 (forecasting); Figure~8.}

\begin{revcomment}{Comment R3.5}
Add sensitivity analysis for fitness function weights.
\end{revcomment}
\rsp{Eleven weight combinations are evaluated over 30 seeds, and the selected weights
are justified by the constraint surface rather than by the headline figure. See
Comment~R2.4.}
\rwhere{Section~6.8 (fitness-weight sensitivity).}

\begin{revcomment}{Comment R3.6}
Include computational cost analysis of CA-WOA.
\end{revcomment}
\rsp{Runtime, peak memory, evaluation counts and convergence behaviour are reported
across the full task-count range and at the smaller counts used in comparable
studies. See Comment~R2.9.}
\rwhere{Section~6.14 (computational cost and convergence).}

\begin{revcomment}{Comment R3.7}
Conduct ablation study on carbon-aware seeding.
\end{revcomment}
\rsp{Reported across six optimisers and two energy models: seeding improves quality
in 59 of 60 cells, and the carbon-blind initialisation control isolates the source
of that benefit. See Comment~R2.3.}
\rwhere{Section~6.7 (ablation); Table~7.}

\begin{revcomment}{Comment R3.8}
Add theoretical convergence analysis.
\end{revcomment}
\rsp{A convergence analysis has been added. The decoded search space is finite, so
the objective takes finitely many values separated by a positive gap, and the
population sequence forms a \emph{non-homogeneous} Markov chain because WOA's
control parameter decreases to zero. Theorem~1 establishes that with elitism the
best-so-far sequence converges almost surely after finitely many improvements.

We also state plainly what does \emph{not} hold: the standard sufficient condition
for convergence to the global optimum fails for WOA, because the reachable set
contracts around the incumbent once exploration switches off. WOA is therefore not
a global-convergence algorithm, and this is a property of WOA rather than of our
extension.

Two consequences justify the design. Corollary~1: precisely because that condition
fails, the limit depends on the initial distribution, so initialisation changes the
answer rather than merely the convergence rate---which is why seeding produces a
statistically significant improvement. Corollary~2: seeding a known schedule bounds
the limit by that schedule's quality, so CA-WOA provably cannot be worse than the
heuristic it seeds. Proposition~1 combines this with Lemma~1 to show that in the
temporal model CA-WOA attains the global optimum in one epoch.}
\rwhere{Section~4.5 (convergence analysis: Theorem~1, Corollaries 1--2, Proposition~1).}

\begin{revcomment}{Comment R3.9}
Include at least one additional carbon region (US or EU).
\end{revcomment}
\rsp{Four additional regions are included: the Republic of Ireland, Northern Ireland,
California and the US Mid-Atlantic. The seeding effect holds in every region
($+0.64$ to $+1.98$ percentage points, all $p < 10^{-9}$).

The Mid-Atlantic region is reported because it is the case where our method does
\emph{not} help. Its carbon intensity varies only $1.78\times$ across the day
against the UK's $12.30\times$, and there CA-WOA loses to the best competing method
by 0.13 points. Temporal shifting requires a low-carbon trough to exploit; where a
grid is uniformly carbon-intense there is nothing to shift into. We regard
identifying this deployment condition as a useful result rather than a negative
one.}
\rwhere{Section~3.1, Table~2; Section~6.12; Figure~6; Section~6.17 (Limitations).}

\begin{revcomment}{Comment R3.10}
Report results with standard deviations for all experiments.
\end{revcomment}
\rsp{All experiments report standard deviations and 95\,\% confidence intervals over
30 seeds, and per-seed values are released. See Comment~R2.6.}
\rwhere{All results tables; archived artefact.}

\begin{revcomment}{Comment R3.11}
Review all citations. The introduction fails to engage with recent, highly relevant
works from 2023--2025. \\
\texttt{10.3390/en17112539}; \texttt{10.3390/electronics13173552};\\
\texttt{10.32604/cmes.2025.065098}; \texttt{10.3390/en17215412}
\end{revcomment}
\rsp{The bibliography has been reviewed in full and expanded from 19 to 31 entries,
every DOI resolved against Crossref. The introduction and related-work sections now
engage with carbon-aware scheduling work from 2022 to 2026, and Table~1 has been
extended accordingly.

On the four suggested references, we have added two and respectfully not the other
two, and we set out our reasoning for the reviewer's consideration. The two
load-forecasting papers are now cited in the forecasting subsection, where a
truthful connection exists: they apply machine learning to forecasting grid-side
quantities, which is the same methodological family as forecasting carbon
intensity. We distinguish the targets explicitly, since carbon intensity depends on
the generation mix rather than on demand alone.

The remaining two concern real-time fault detection and isolation in power systems,
and the assessment of power-transformer oil deterioration. We were unable to
identify a substantive connection to cloud task scheduling, carbon-aware computing
or metaheuristic optimisation that would let us cite them accurately, and we
preferred not to include citations we could not motivate in the text. We would of
course reconsider if the reviewer can indicate the intended connection.}
\rwhere{Section~2.4; Table~1; Bibliography.}

\bigskip
\hrule
\bigskip

\section*{Items awaiting supervisor input}
The following are left blank deliberately, as agreed.

\begin{revcomment}{Item S.1}
Whether a full MILP solve should be attempted for the camera-ready version to
provide optimality gaps, or whether the theoretical treatment in Section~4.4 is
sufficient as submitted.
\end{revcomment}
\rpending

\begin{revcomment}{Item S.2}
Whether the response to Comment~R3.11 regarding the two uncited suggested
references should be softened, expanded, or removed.
\end{revcomment}
\rpending

\begin{revcomment}{Item S.3}
Whether any declaration regarding the use of generative AI tools during the
revision is required by the submission venue, and if so its wording.
\end{revcomment}
\rpending

\end{document}
